\section{Chapter Prerequisites}\label{limits_prereq}\thispagestyle{fancy}

\prereqIntro

\subsection{Functions}

A \textbf{function} $f$ is a rule that assigns each element $x$ from a set (called the domain) to exactly one element, called $f(x)$, in another set. Unless we say otherwise, the \textbf{domain} is the set of all real numbers for which the rule makes sense and defines a real number. All possible values of $f(x)$ are called the \textbf{range} of $f$. We use four ways to represent a function.
\begin{itemize}
\begin{multicols}{2}\lxAddClass{columns2}
\item By a graph
\item By an explicit formula
\item By a table of values
\item By a verbal description
\end{multicols}
\end{itemize}

Throughout the book we will use several representations of any given function to help give us a better understanding of the problem. The graphs in \autoref{prereq_basic_graphs} contain most of the base functions we can use to build other functions using transformations.

\noindent%
\flushinner{%
\newcommand{\picturescale}{.28}
\begin{minipage}[t]{1.27\linewidth}\tagstructbegin{tag=Sect}\noindent%
\captionsetup{type=figure}%
\tagpdfsetup{table/header-rows={2,4}}
\begin{tabular}{c c c c c}
\begin{tikzpicture}[alt={A U shape, touching the origin.},scale = \picturescale]
\begin{axis}[axis y line=middle, axis x line=middle, xmin=-2.2, xmax=2.2, ymin=-.5, ymax=2.2,name=myplot, ytick={0,1,2},axis equal]
\addplot[draw={\colorone}, domain=-2:2, thick, smooth]{x*x};
\end{axis}
\node [right] at (myplot.right of origin) {\small $x$};
\node [above] at (myplot.above origin) {\small $y$};
\end{tikzpicture}
&
\begin{tikzpicture}[alt={Flat at the origin, going up to the right and down to the left.},scale = \picturescale]
\begin{axis}[axis y line=middle, axis x line=middle, xmin=-2.2, xmax=2.2, ymin=-2.2, ymax=2.2,name=myplot,axis equal]
\addplot[draw={\colorone}, domain=-2:2, thick, smooth]{x*x*x};
\end{axis}
\node [right] at (myplot.right of origin) {\small $x$};
\node [above] at (myplot.above origin) {\small $y$};
\end{tikzpicture}
&
\begin{tikzpicture}[alt={Starting along the x axis, and then swooping upward.},scale = \picturescale]
\begin{axis}[axis y line=middle, axis x line=middle, xmin=-2.2, xmax=2.2, ymin=-.5, ymax=3.2,name=myplot, ytick={0, 1,2,3},axis equal]
\addplot[draw={\colorone}, domain=-2:2, thick, smooth]{exp(x)};
\end{axis}
\node [right] at (myplot.right of origin) {\small $x$};
\node [above] at (myplot.above origin) {\small $y$};
\end{tikzpicture}
&
\begin{tikzpicture}[alt={A wave passing through the origin.},scale = \picturescale]
\begin{axis}[axis y line=middle, axis x line=middle, xmin=-.5, xmax=6.5, ymin=-1.2, ymax=1.2,name=myplot, ytick={-1,0,1}, xtick={-6.28318, -4.7123889, -3.14159, -1.5708, 1.5708, 3.14159, 4.7123889, 6.28318}, xticklabels={-$2\pi$, $-3\pi/2$,$-\pi$, $-\pi/2$, $\pi/2$,$\pi$, $3\pi/2$, $2\pi$},axis equal]
\addplot[draw={\colorone}, domain=-.5:6.4, thick, smooth]{sin deg(x)};
\end{axis}
\node [right] at (myplot.right of origin) {\small $x$};
\node [above] at (myplot.above origin) {\small $y$};
\end{tikzpicture}
&
\begin{tikzpicture}[alt={A V touching the origin.},scale = \picturescale]
\begin{axis}[axis y line=middle, axis x line=middle, xmin=-2.2, xmax=2.2, ymin=-.5, ymax=2.2,name=myplot, ytick={0,1,2},axis equal]
\addplot[draw={\colorone}, domain=-2:2, thick, smooth]{abs x};
\end{axis}
\node [right] at (myplot.right of origin) {\small $x$};
\node [above] at (myplot.above origin) {\small $y$};
\end{tikzpicture}
\\ $y=x^2$ & $y=x^3$ & $y=e^x$ & $y=\sin x$ & $y=\abs{x}$
 \\
\begin{tikzpicture}[alt={Starting at the origin, moving upward and then leveling off.},scale = \picturescale]
\begin{axis}[axis y line=middle, axis x line=middle, xmin=-.5, xmax=3.2, ymin=-.5, ymax=2.2,name=myplot, ytick={0,1,2}, xtick={0,1,2,3},axis equal]
\addplot[draw={\colorone}, domain=-0:3, thick, smooth,samples=100]{sqrt x};
\end{axis}
\node [right] at (myplot.right of origin) {\small $x$};
\node [above] at (myplot.above origin) {\small $y$};
\end{tikzpicture}
&
\begin{tikzpicture}[alt={Vertical at the origin, moving up to the right and down to the left, but leveling off as it gets away from the origin.},scale = \picturescale]
\begin{axis}[axis y line=middle, axis x line=middle, xmin=-2.2, xmax=2.2, ymin=-2.2, ymax=2.2,name=myplot,axis equal]
\addplot[draw={\colorone}, domain=-2:2, thick, smooth,samples=100]({x*x*x},{x});
\end{axis}
\node [right] at (myplot.right of origin) {\small $x$};
\node [above] at (myplot.above origin) {\small $y$};
\end{tikzpicture}
&
\begin{tikzpicture}[alt={Starting vertically next to the negative y axis, moving up and to the right as it levels off.},scale = \picturescale]
\begin{axis}[axis y line=middle, axis x line=middle, xmin=-.5, xmax=3.2, ymin=-2.2, ymax=2.2,name=myplot, ytick={-2,-1,0, 1,2},xtick={0,1,2,3},axis equal]
\addplot[draw={\colorone}, domain=.01:3, thick, smooth]{ln x};
\end{axis}
\node [right] at (myplot.right of origin) {\small $x$};
\node [above] at (myplot.above origin) {\small $y$};
\end{tikzpicture}
&
\begin{tikzpicture}[alt={A wave cresting above the origin.},scale = \picturescale]
\begin{axis}[axis y line=middle, axis x line=middle, xmin=-.5, xmax=6.5, ymin=-1.2, ymax=1.2,name=myplot, ytick={-1,0,1}, xtick={-6.28318, -4.7123889, -3.14159, -1.5708, 1.5708, 3.14159, 4.7123889, 6.28318}, xticklabels={-$2\pi$, $-3\pi/2$,$-\pi$, $-\pi/2$, $\pi/2$,$\pi$, $3\pi/2$, $2\pi$},axis equal]
\addplot[draw={\colorone}, domain=-.5:6.4, thick, smooth]{cos deg(x)};
\end{axis}
\node [right] at (myplot.right of origin) {\small $x$};
\node [above] at (myplot.above origin) {\small $y$};
\end{tikzpicture}
&
\begin{tikzpicture}[alt={Starting just below the negative x axis, moving toward the negative y axis.  Then the graph jumps to the positive y axis and moves toward the positive x axis.},scale = \picturescale]
\begin{axis}[axis y line=middle, axis x line=middle, xmin=-3.2, xmax=3.2, ymin=-3.2, ymax=3.2,name=myplot, ytick={-3,-2,-1,...,3},axis equal,xtick={2}]
\addplot[draw={\colorone}, domain=-3:-0.2, thick, smooth]{1/x};
\addplot[draw={\colorone}, domain=0.2:3, thick, smooth]{1/x};
\end{axis}
\node [right] at (myplot.right of origin) {\small $x$};
\node [above] at (myplot.above origin) {\small $y$};
\end{tikzpicture}\\
$y=\sqrt x$ & $y=\sqrt[3]x$ & $y=\ln x$ & $y=\cos x$ & $y=\dfrac{1}{x}$
\end{tabular}
\caption{Basic Function Graphs}\label{prereq_basic_graphs}\tagstructend
\end{minipage}
}%

We will often transform these functions into other functions as given in the next two figures.\\
\noindent\begin{minipage}[t]{\linewidth}\tagstructbegin{tag=Sect}\noindent%
\captionsetup{type=figure}%
\centering
\tagpdfsetup{table/header-rows={1}}
\begin{tabular}{l c}\lxBeginTableHead
The function & shifts $f(x)$\\\lxEndTableHead\midrule
$y=f(x)+c$ & $c$ units upward\\
$y=f(x)-c$ & $c$ units downward\\
$y=f(x+c)$ & $c$ units left\\
$y=f(x-c)$ & $c$ units right\\
\end{tabular}
\caption{Translations of Basic Functions with $c>0$}\tagstructend
\end{minipage}\bigskip

\leavevmode
\noindent\begin{minipage}[t]{\linewidth}\tagstructbegin{tag=Sect}\noindent%
\captionsetup{type=figure}%
\centering
\tagpdfsetup{table/header-rows={1}}
\begin{tabular}{l c}\lxBeginTableHead
The function &  transforms $f(x)$ by\\\lxEndTableHead\midrule
$y=cf(x)$ & stretching vertically by a factor of $c$\\
$y=\frac{1}{c} f(x)$ & shrinking vertically by a factor of $c$\\
$y=f(cx)$ & shrinking horizontally by a factor of $c$\\
$y=f(\frac{x}{c})$ & stretching horizontally by a factor of $c$\\
$y=-f(x)$ & reflecting about the $x$-axis\\
$y=f(-x)$ & reflecting about the $y$-axis\\
\end{tabular}
\caption{Scaling Basic Functions with $c>1$}\tagstructend
\end{minipage}

%\begin{example}[Sketching with basic transformations]\label{ex_prereq_sketch}%
%Sketch the graph of the following functions using the base function and the appropriate transformations.
%\begin{multicols}{2}
%\begin{enumerate}\lxAddClass{columns2}
%\item $y=\displaystyle \frac{1}{x+3}$
%\item $y=\sqrt{x+3}+1$
%\item $y=\abs{x-4}$
%\item $y=3\cos x+2$
%\item $y=4\abs{x}+1$
%\item $y=-\frac{1}{3}(x-2)^2+3$
%\item $y=(x-3)^3$
%\item$y=\abs{\sin 2x}$
%\end{multicols}
%\end{enumerate}
%\solution
%\flushinner{%
%\tagpdfsetup{table/tagging=presentation}
%\begin{tabular}{ c c c c }
%\begin{tikzpicture}[alt={Graph starting at (-4,-1) then going down toward (-3,-infty), jumping to (-3,infty), moving downward to curve right and finish near (4,0.14).},scale = .5]
%\begin{axis}[axis y line=middle, axis x line=middle, xmin=-4, xmax=4, ymin=-4, ymax=4,name=myplot, ytick={0, 1,2}, xtick={0,1,2,3}]
%\addplot[draw={\colorone}, domain=-3:4, thick, smooth,samples=100]{1/(x+3)};
%\addplot[draw={\colorone}, domain=-4:-3, thick, smooth,samples=100]{1/(x+3)};
%\end{axis}
%\node [right] at (myplot.right of origin) {\scriptsize $x$};
%\node [above] at (myplot.above origin) {\scriptsize $y$};
%\end{tikzpicture}
%&
%\begin{tikzpicture}[alt={Graph starting at (-3,1), moving upward and curving to the right to finish near (4,3.6).},scale = .5]
%\begin{axis}[axis y line=middle, axis x line=middle, xmin=-4, xmax=4, ymin=-4, ymax=4,name=myplot, ytick={0, 1,2}, xtick={0,1,2,3}]
%\addplot[draw={\colorone}, domain=-3:4, thick, smooth,samples=100]{sqrt(x+3)+1};
%\end{axis}
%\node [right] at (myplot.right of origin) {\scriptsize $x$};
%\node [above] at (myplot.above origin) {\scriptsize $y$};
%\end{tikzpicture}
%&
%\begin{tikzpicture}[alt={Graph starting at (-1,5), moving straight to (4,0), making a right angle turn upward to finish at (5,1).},scale = .5]
%\begin{axis}[axis y line=middle, axis x line=middle, xmin=-5, xmax=5, ymin=-5, ymax=5,name=myplot, ytick={0, 1,2}, xtick={0,1,2,3}]
%\addplot[draw={\colorone}, domain=-5:5, thick, smooth,samples=100]{abs(x-4)};
%\end{axis}
%\node [right] at (myplot.right of origin) {\scriptsize $x$};
%\node [above] at (myplot.above origin) {\scriptsize $y$};
%\end{tikzpicture}
%\end{tabular}}
%end{example}

\subsection{Domain}

We said above that domain is the set of real numbers for which the function (rule) defines a real number and makes sense. Ask yourself, ``what values can I put into the function and get a real value out?'' There are generally two key expressions that will limit the domain of a function from all real numbers. We may not divide by zero and we may not have a negative number underneath an even root. The following examples illustrate how we restrict the domain when we see these expressions.

\begin{example}[Finding a domain]\label{prereq_domain_1}%
Find the domain of the function $f(x)=\sqrt{x-4}$.
\solution
The square root of a negative number is not defined as a real number so the domain of $f$ will be all real numbers for which $x-4\geq 0$ which is $x \geq 4$. In interval notation, this is $[4,\infty)$.
\end{example}

\begin{example}[Finding a domain]\label{prereq_domain_2}%
Find the domain of the function $g(x)= \dfrac{3}{x^2-9}$.
\solution
We cannot divide by zero so we factor the denominator of $g$ and exclude those values where the denominator is zero.
\[g(x)=\frac3{x^2-9}=\frac3{(x-3)(x+3)}\]
We see that $x\neq 3,-3$ for $g$ to be defined, which is written in interval notation as $(-\infty,-3)\cup (-3,3)\cup (3,\infty)$.
\end{example}

\begin{example}[Finding a domain]\label{prereq_domain_3}%
Find the domain of the function $h(x)=\dfrac1{\sqrt{x^2-4}}$
\solution
For $h$ to be defined as a real number we must have $x^2-4>0$. This is equivalent to $(x-2)(x+2)>0$ and we create a sign chart:

\begin{center}
\includecodegraphics{figures/matrices/prereqMatrix}
\end{center}

This shows that $x^2-4$ will be greater than zero on $(-\infty,-2)\cup(2,\infty)$.
\end{example}

% todo write a trigonometry review

\printexercises{exercises/01-00-exercises}
